% Fragmento conservado; véase MANIFIESTO_ANTECEDENTES_LECTURA.json.
\section{Registre dodécaphasique et constante rationnelle de phase}
\label{sec:k-registro}

La génération des coordonnées régionales n'épuise pas l'information de l'état APP--TRIT--TPK. Une lecture archimédienne détermine une valeur au moyen de cylindres compatibles ; une autre lecture conserve l'opposition des feuillets, l'orientation de la route et le registre des incidences traversées. Le registre dodécaphasique appartient à cette seconde lecture avant d'intervenir dans la clôture d'alpha. Il n'est donc pas défini en soustrayant une constante physique à une combinaison de valeurs préalablement connues. Il est reconstruit à partir d'une coordonnée signée du même état arithmétique avec mémoire de trajectoire qui produit les coordonnées régionales.

La distinction est nécessaire même lorsque deux objets possèdent douze
coordonnées. Un mot de douze blocs du catalogue, une liste de douze
événements orientés, un vecteur entier signé et un registre dodécaphasique en base mille sont des
objets distincts. Leur comparaison requiert les applications qui les
relient. L'égalité de longueur ne fournit pas ces applications, et la
perte des signes n'est pas une simple simplification de notation.

Ce chapitre maintient la délimitation causale fixée dans les chapitres précédents :
les neuf positions APP des axes horizontal et vertical produisent les
feuillets arithmétiques ; TRIT type le régime et l'orientation ; TPK sélectionne,
transporte, met à jour et prolonge, en conservant le résidu, le quotient, la retenue,
le feuillet, la route, la frontière et la mémoire. À partir de la structure discrète conjointe du
continu s'obtiennent les coordonnées utilisées ici. La chaîne
\(w_6\to w_{12}\to w_{18}\to w_{24}\to w_{30}\to R_{36}\to G_9\)
n'est pas remplacée par un calendrier fini. L'ordre
\(U_t\to\Gamma_9\to K_{\mathrm{ph}}\) n'est pas non plus inversé : la dernière application est une
lecture à mémoire réduite, et non le générateur du registre signé.

% Fragmento conservado; véase MANIFIESTO_ANTECEDENTES_LECTURA.json.
% Adición propuesta al inicio de registro_k.tex, después de su introducción.
% Conserva todas las definiciones, fórmulas y demostraciones actuales.
% Procedencia: registro_k.tex, §§k-hadamard/k-transversal/k-kappas;
% k_reversibilidad.tex; incidencia_articulacion.tex, inc-ampl-cuaterna.
\paragraph{Coordonnées réversibles et détermination incidencielle.}
\label{ampl-alfa-k-intro-reversible}
Le résultat qui organise ce chapitre est l'équivalence entre trois
systèmes de coordonnées d'un registre préalablement produit par
APP--TRIT--TPK. La transformation de Hadamard relie le registre
signé au registre entier ; l'extracteur de différences relie
ce dernier à ses variations transversales et à sa composante uniforme :
\[
 U\underset{\mathcal H_{12}}{\overset{\mathcal H_{12}/4}{\longleftrightarrow}}K
 \overset{(D_3,D_4,Q)}{\longleftrightarrow}
 (D_3K,D_4K,Q(K)).
\]
La première équivalence s'établit sur le sous-réseau entier
$\mathcal L_H$ ; la seconde, sur l'image de l'extracteur. Une fois fixés la
base $1000$, les $12$ positions et l'origine de phase, l'évaluation
$K\mapsto\kappa_{\mathrm{per}}$ ajoute un quatrième système de coordonnées
réversible : la multiplication par $1000^{12}-1$ suivie des divisions
euclidiennes restitue le registre complet. Les démonstrations de
\S\ref{subsec:k-hadamard}, \S\ref{subsec:k-transversal} et
\S\ref{subsec:k-kappas} établissent ces équivalences dans leurs domaines
respectifs.

Deux compositions agissent ensuite sur ce registre. La composition
positionnelle combine ses composantes avec les coordonnées de clôture,
de propagation et d'autoéchelle, et sa normalisation détermine $\alpha_A$. La
composition incidencielle applique le seuil cubique et obtient
\[
 K\longmapsto B_K=\{j:K_j\geq729\}=\{4,6,9,10\}.
\]
Le registre intégral permet de reconstruire les deux compositions. Le support
$B_K$ conserve spécifiquement l'appartenance à la classe sélectionnée
par le seuil. L'identité $B_K=H_e\cap P_\pi$ de
(8.12) de l'Article I (provenance du prolongement exceptionnel, conservée dans le registre technique) relie cette extraction entière à l'incidence
des codes régionaux. Le support négatif de son bilan avec $K$
complète le drapeau $B_K\subset H^-_\alpha$ développé dans le
prolongement exceptionnel. Ainsi, la fonction de $K$ est définie par
une transformation inversible et par des applications ultérieures
précises ; l'information supplémentaire de profondeur et de transport demeure
dans l'état dont procède le registre.



\subsection{Extension cyclique et graduation de mémoire}
\label{subsec:k-calendario}

Le calendrier observable contient cent huit positions, organisées en douze fenêtres nonadiques. Avec des indices positifs, soient
\[
 E_m=\{9(m-1)+1,\ldots,9m\},\qquad 1\leq m\leq12.
\]
Le support ordonné \(E_1\sqcup\cdots\sqcup E_{12}\) conserve la position et l'
appartenance à une fenêtre. Le groupe \(C_{108}=\mathbb Z/108\mathbb Z\)
conserve, en revanche, la loi de translation cyclique. Pour le représentant
\(0\leq\widetilde\theta<108\), les coordonnées
\[
 \beta(\theta)=\left[\left\lfloor\frac{\widetilde\theta}{9}\right\rfloor\right]_{12},
 \qquad r(\theta)=1+(\widetilde\theta\bmod9)
\]
séparent le secteur dodécaphasique et la position nonadique intérieure. Aucune d'elles ne conserve à elle seule le nombre total de tours.

\begin{proposition}[La retenue du calendrier]
\label{prop:k-extension}
Les applications
\[
 \iota:C_{12}\longrightarrow C_{108},\quad[b]\longmapsto[9b],
 \qquad p:C_{108}\longrightarrow C_9,\quad[\theta]\longmapsto[\theta]_9
\]
forment une suite exacte non scindée
\[
 0\longrightarrow C_{12}\xrightarrow{\iota}C_{108}
 \xrightarrow{p}C_9\longrightarrow0.
\]
Pour la section ensembliste qui choisit les représentants \(0,\ldots,8\), le défaut d'additivité est la retenue
\[
 c(r,r')=\left[\left\lfloor\frac{\widetilde r+\widetilde r'}9\right\rfloor\right]_{12}.
\]
\end{proposition}

\begin{proof}
Le noyau de \(p\) est formé des multiples de neuf et coïncide avec
l'image de \(\iota\). La réduction est surjective et la multiplication
par neuf est injective sur le domaine indiqué. Si la suite était
scindée, le groupe central serait isomorphe à \(C_{12}\times C_9\), dont l'
exposant est \(\operatorname{mcm}(12,9)=36\) ; l'exposant de
\(C_{108}\) est 108. La contradiction prouve la non-scission. Enfin,
la division euclidienne de \(\widetilde r+\widetilde r'\) par neuf donne la
retenue écrite et vérifie l'identité de la section.
\end{proof}

Le résultat explique une différence qui réapparaîtra dans alpha : revenir à la
phase n'équivaut pas à revenir à l'état. Si \(q_{\mathrm{mem}}\) compte les
tours nonadiques, le calendrier ne conserve que
\(\beta=[q_{\mathrm{mem}}]_{12}\). Après cent huit pas,
\(\beta\) revient, mais \(q_{\mathrm{mem}}\) augmente de douze. Le
registre conserve en outre le cylindre, le livre d'incidences et la mémoire
verticale. La clôture d'une coordinatisation finie ne transforme pas cet historique en un
cycle sans mémoire.

\subsection{Le registre orienté et la coordinatisation intégrale \texorpdfstring{\(1+5+6\)}{1+5+6}}
\label{subsec:k-registro-integral}

Pour une histoire produite par TPK, soient \(\tau_m\) la sous-route sur
\(E_m\), \(\sigma_m\) son orientation, \(\kappa_m\) la retenue de
frontière et \(\Lambda_m\) le livre local des incidences. Le registre est
\[
 \mathcal R_{12}(x)=
 \bigl((E_m,\tau_m,\sigma_m,\kappa_m,\Lambda_m)\bigr)_{m=1}^{12}.
\]
Son domaine est l'espace des histoires arithmétiques avec mémoire de trajectoire, et non l'ensemble des
étiquettes de phase. La projection sur le calendrier oublie précisément certaines
des données dont le lecteur signé a besoin.

Une coordinatisation initiale du registre peut s'écrire au moyen d'événements orientés. La construction de référence~\textup{(provenance conservée dans le registre technique)} fixe l'ensemble des codes d'événement \(\{2,4,6,8\}\) et l'orientation sectorielle
\[
 \Sigma_{12}=(+,+,+,-,-,-,+,+,+,-,-,-).
\]
Si \(d_t\) est le code de direction au pas \(t\), on obtient
\[
 e_i=\Sigma_{12}(i+1)
 \#\{t\in E_{i+1}:d_t\in\{2,4,6,8\}\},
 \qquad0\leq i<12.
\]
La formule maintient distincts le recensement non négatif des événements et
l'orientation. Elle n'identifie pas une direction cardinale à un signe TRIT. Le
vecteur \(e\), dont les composantes sont des recensements locaux signés, n'est pas non plus
encore le vecteur \(U\) qui apparaîtra plus bas.

Les positions opposées déterminent, pour \(0\leq i<6\),
\[
 p_i=e_i+e_{i+6},\qquad a_i=e_i-e_{i+6}.
\]
Avec
\[
 s_C=\sum_{i=0}^{5}p_i,\qquad M_i=p_i-p_5\quad(0\leq i<5),
\]
on définit la coordinatisation intégrale
\[
 \mathcal L_{12}(e)
 =(s_C,M_0,\ldots,M_4,a_0,\ldots,a_5).
\]
La décomposition \(1+5+6\) n'est pas une répartition arbitraire des coordonnées :
une coordonnée conserve la charge, cinq comparent les paires par rapport à une
paire de référence et six conservent l'opposition des demi-couronnes.

\begin{lemma}[Inversion de la coordinatisation intégrale]
\label{lem:k-carta-integral}
Un tuple entier \((s_C,M_0,\ldots,M_4,a_0,\ldots,a_5)\) appartient à l'image de \(\mathcal L_{12}\) si et seulement si
\[
 s_C-\sum_{i=0}^{4}M_i\equiv0\pmod6,
 \qquad p_i\equiv a_i\pmod2\quad(0\leq i<6),
\]
où
\[
 p_5=\frac{s_C-\sum_{i=0}^{4}M_i}{6},\qquad p_i=p_5+M_i\quad(i<5).
\]
Dans ce cas, sa préimage est unique et est donnée par
\[
 e_i=\frac{p_i+a_i}{2},\qquad e_{i+6}=\frac{p_i-a_i}{2}.
\]
\end{lemma}

\begin{proof}
La somme des cinq marges est \(s_C-6p_5\), d'où l'on obtient la
division par six. Une fois les \(p_i\) reconstruits, chaque paire d'équations
\(p_i=e_i+e_{i+6}\), \(a_i=e_i-e_{i+6}\) possède la solution indiquée.
Cette solution est entière exactement sous la congruence de parité. Toutes
les opérations s'inversent, de sorte qu'aucune liberté résiduelle ne subsiste.
\end{proof}

Cette réversibilité montre quelle information est conservée ; elle n'autorise pas à
supprimer le reste du livre. Le registre complet possède un domaine plus vaste
que cette coordinatisation des recensements. L'extraction des canaux qui alimentent
\(U\) utilise ce registre complet et ne se déduit pas de la liste
\((e_i)\) par une identification de noms.

\subsection{Deux lectures de l'état terminal}
\label{subsec:k-terminal}

Désignons par \(x_{\mathrm{term}}\) l'état produit par la composition de sélection, de transport et de mise à jour jusqu'à la coupe terminale. L'écriture opératorielle est
\[
 x_{\mathrm{term}}=
 (\mathcal U_{108}\circ\cdots\circ\mathcal U_1)(x_{\mathrm{seed}}),
 \qquad \mathcal U_t=\mathsf{Upd}_t\circ\mathsf{Tra}_t\circ\mathsf{Sel}_t.
\]
Ses deux lectures pertinentes sont
\[
 x_{\mathrm{term}}
 \xrightarrow{(\pi_{\mathrm{term}},R_{\mathrm{sgn}})}
 (B_{\mathrm{term}},U).
\]
La matrice \(B_{\mathrm{term}}\) retient la lecture ternaire visible ; le vecteur \(U\) conserve une autre coordonnée de la même histoire. On n'insère pas entre eux une flèche \(B_{\mathrm{term}}\to U\).

Le sélecteur terminal de la construction de référence~\textup{(provenance conservée dans le registre technique)} opère sur des catalogues de tailles 196, 197 et 196. Les signatures locales réduisent leurs \(196\cdot197\cdot196=7\,567\,952\) triplets à 331, et la marge des colonnes \(q_{\partial}=(1,2,2,1,2,0)\) laisse deux matrices :
\[
 B_0=\begin{pmatrix}0&2&0&2&2&0\\0&0&1&2&2&0\\1&0&1&0&1&0\end{pmatrix},
 \qquad
 B_1=\begin{pmatrix}0&2&1&0&2&0\\0&0&1&2&2&0\\1&0&0&2&1&0\end{pmatrix}.
\]
L'orientation des lignes est \(\sigma=(1,1,-1)=(1,1,2)\) dans
\(\mathbb F_3^3\). Les charges des lignes des candidates sont
\((0,2,0)\) et \((2,2,1)\). Seule la seconde appartient à
\(\langle\sigma\rangle=\{(0,0,0),(1,1,2),(2,2,1)\}\) ; par conséquent,
\(B_{\mathrm{term}}=B_1\). Cette vérification prouve la dernière étape du
sélecteur à partir du recensement déclaré. Elle ne transforme pas la marge des colonnes
en une conséquence de la charge des lignes, et ne remplace pas l'énumération
précédente par l'inspection de ces deux matrices.

Le lecteur signé possède la composition typée
\begin{equation}
 R_{\mathrm{sgn}}
 =\Pi_H\circ\mathcal E_{108}^{90,120}\circ\mathcal R_{12}:
 \mathcal X_{\mathrm{TPK}}^{\mathrm{term,mem}}
 \longrightarrow\mathbb Z^{12}.
 \label{eq:k-lector-firmado}
\end{equation}
Le registre des canaux détermine
\[
 \mathcal E_{108}^{90,120}(\mathcal R_{12}(x_{\mathrm{term}}))
 =(b^{90},b^{120},Q_{\mathrm{TPK}}),
\]
avec les coordonnées
\begin{align*}
 b^{90}={}&(-495,-116,-684,108,601,-90,\\
            &\hspace{19mm}-173,-88,313,560,-397,461),\\
 b^{120}={}&(-425,-281,-481,671,-255,30,\\
             &\hspace{19mm}475,-543,680,251,6,-128),\\
 Q_{\mathrm{TPK}}={}&6263.
\end{align*}
Ces quantités sont utilisées comme coordonnées du registre préalable, non
comme paramètres ajustés au moyen d’alpha. La reconstruction démontrée
dans les paragraphes suivants commence exactement à cette sortie typée.
La trace matérielle de l’extracteur antérieur et la portée de sa transcription
autonome sont conservées dans la note technique de provenance ; l’arithmétique
ultérieure n’est pas présentée comme un substitut à cette trace.

\subsection{Du registre des canaux au vecteur signé}
\label{subsec:k-pi-h}

Le lecteur harmonique \(\Pi_H\) peut s'écrire sans utiliser ni \(K\) ni
alpha. Pour éviter de confondre ses sommes orbitales avec les blocs d'alpha,
nous les noterons \(s_1,s_2,s_3\). Soient
\[
 B_r=\sum_{j=0}^{3}b^{120}_{r+3j},\qquad
 d_{r,j}=b^{90}_{r+3j},\qquad1\leq r\leq3,
\]
et
\begin{equation}
 s_1=\frac{Q_{\mathrm{TPK}}+2B_1+B_2}{3},\qquad
 s_2=s_1-B_1,\qquad s_3=s_2-B_2.
 \label{eq:k-sumas-armonicas}
\end{equation}
La division par trois appartient au lecteur des trois orbites. Son caractère entier est vérifié sur les sorties compatibles du registre ; il n'est pas supposé pour un élément arbitraire de \(\mathbb Z^{25}\).

Les trois blocs signés sont
\begin{equation}
 U^{(r)}=(s_r,d_{r,1}-d_{r,3},d_{r,0}+d_{r,2},d_{r,0}-d_{r,2}).
 \label{eq:k-bloques-firmados}
\end{equation}
La substitution entière donne
\[
 (B_1,B_2,B_3)=(972,-1073,101),\qquad
 (s_1,s_2,s_3)=(2378,1406,2479),
\]
et, par conséquent,
\begin{align*}
 U^{(1)}&=(2378,-452,-668,-322),\\
 U^{(2)}&=(1406,998,-204,-28),\\
 U^{(3)}&=(2479,-551,-371,-997).
\end{align*}
En entrelaçant les colonnes, on obtient
\begin{equation}
 U=(2378,1406,2479,-452,998,-551,
       -668,-204,-371,-322,-28,-997).
 \label{eq:k-u-firmado}
\end{equation}
La présence de valeurs négatives et de composantes supérieures à 999
montre que \(U\) n'est pas un mot de chiffres en base mille. Ce n'est pas non plus
\(U_6\Vert U_6\), concaténation du catalogue de période six. La différence
concerne le domaine et l'information conservée, et non seulement les valeurs.

L'ordre constructif jusqu'ici est
\[
 x_{\mathrm{term}}\longrightarrow\mathcal R_{12}
 \longrightarrow(b^{90},b^{120},Q_{\mathrm{TPK}})
 \longrightarrow U.
\]
L'inversion qui suit produit le registre dodécaphasique. Le fait que nous puissions ensuite récupérer les canaux à partir du registre dodécaphasique constitue un contrôle de réversibilité ; il n'inverse pas cet ordre causal.

\subsection{Hadamard par orbites et reconstruction entière}
\label{subsec:k-hadamard}

Dans le cycle à douze positions, le pas trois détermine trois orbites :
\[
 (1,4,7,10),\qquad(2,5,8,11),\qquad(3,6,9,12).
\]
La séparation de quatre secteurs de chaque orbite est celle qui, une fois
l'origine et l'orientation fixées, admet la lecture angulaire de quatre-vingt-dix degrés. La
transformation n'agit pas sur trois quadruplets contigus de l'ordre naturel.

Soit
\[
 H_4=\begin{pmatrix}
 1&1&1&1\\1&1&-1&-1\\1&-1&1&-1\\1&-1&-1&1
 \end{pmatrix}.
\]
Si \(P_{\mathrm{orb}}\) réordonne selon les trois orbites indiquées, nous définissons
\[
 \mathcal H_{12}=P_{\mathrm{orb}}^{-1}
 (H_4\oplus H_4\oplus H_4)P_{\mathrm{orb}}.
\]

\begin{lemma}[Inversion de Hadamard et intégralité]
\label{lem:k-hadamard}
On a
\[
 H_4^2=4I_4,\qquad H_4^{-1}=\frac14H_4,\qquad\det H_4=-16.
\]
Pour \(u\in\mathbb Z^4\), la préimage \(H_4^{-1}u\) est entière si et
seulement si \(H_4u\equiv0\pmod4\).
\end{lemma}

\begin{proof}
Les lignes sont orthogonales et le carré de leur norme vaut quatre ; la matrice est symétrique. Cela donne \(H_4^2=4I_4\) et l'inverse. Une élimination entière, ou le développement du déterminant, donne le signe négatif et la valeur \(-16\). La condition de caractère entier est exactement la divisibilité par quatre des quatre coordonnées de \(H_4u\).
\end{proof}

\begin{definition}[Transformation entière dodécaphasique]
Le sous-réseau des registres signés et sa coordinatisation entière sont
\[
 \mathcal L_H=\{u\in\ZZ^{12}:\mathcal H_{12}u\equiv0\pmod4\},
 \qquad \mathscr K(u)=\tfrac14\mathcal H_{12}u.
\]
L'application \(\mathscr K:\mathcal L_H\to\ZZ^{12}\) est un isomorphisme
de groupes abéliens, d'inverse \(k\mapsto\mathcal H_{12}k\).
Le registre canonique \(K\) est l'image du registre signé généré
par les opérations précédentes, avec une origine de phase et une orientation fixées.
\end{definition}

\begin{proof}
La congruence définit exactement l'intégralité de \(\mathscr K(u)\).
L'identité \(\mathcal H_{12}^2=4I\) démontre les deux compositions
inverses. En particulier, tout \(k\in\ZZ^{12}\) a pour préimage
\(\mathcal H_{12}k\in\mathcal L_H\).
\end{proof}

\begin{proposition}[Registre entier canonique]
\label{prop:k-sello}
L'inversion \(K^{(r)}=H_4U^{(r)}/4\) des blocs de
\eqref{eq:k-bloques-firmados} produit
\begin{align*}
 K^{(1)}&=(234,729,621,794),\\
 K^{(2)}&=(543,659,58,146),\\
 K^{(3)}&=(140,824,914,601).
\end{align*}
Dans l'ordre naturel, le registre dodécaphasique est
\begin{equation}
 K=(234,543,140,729,659,824,621,058,914,794,146,601).
 \label{eq:k-vector}
\end{equation}
C'est l'unique préimage rationnelle de \(U\), et elle appartient à \(\{0,\ldots,999\}^{12}\).
\end{proposition}

\begin{proof}
Les trois produits avant division sont
\begin{align*}
 H_4U^{(1)}&=(936,2916,2484,3176),\\
 H_4U^{(2)}&=(2172,2636,232,584),\\
 H_4U^{(3)}&=(560,3296,3656,2404).
\end{align*}
Toutes les coordonnées sont divisibles par quatre. Diviser et défaire la
permutation des orbites donne \eqref{eq:k-vector}. L'unicité procède de l'
inversibilité rationnelle, et non d'une recherche parmi des candidats proches d'
une valeur physique.
\end{proof}

Le caractère entier peut également se formuler en termes de réseaux. Les diviseurs déterminantaux de \(H_4\) sont \(1,2,4,16\) : les entrées ont pour plus grand commun diviseur un ; les mineurs d'ordre deux sont pairs et l'un d'eux est de module deux ; ceux d'ordre trois sont des multiples de quatre et l'un d'eux est de module quatre. On en déduit
\[
 \operatorname{SNF}(H_4)=\operatorname{diag}(1,2,2,4),
\]
et pour la transformation globale,
\[
 \operatorname{coker}\mathcal H_{12}
 \cong(\mathbb Z/2\mathbb Z)^6\oplus(\mathbb Z/4\mathbb Z)^3.
\]
L'indice de l'image est \(16^3=4096\). Ainsi, tout vecteur signé
entier n'admet pas nécessairement une préimage entière ; le vecteur obtenu satisfait les
congruences précises exigées par l'inversion. Le facteur \(1/4\) est
imposé par la matrice et n'a pas le statut d'un coefficient libre.

\subsection{Différences cycliques et composante uniforme}
\label{subsec:k-transversal}

Avec des indices cycliques modulo douze, soient
\[
 (D_3z)_m=z_m-z_{m+3},\qquad(D_4z)_m=z_m-z_{m+4},
 \qquad Q(z)=\sum_{m=1}^{12}z_m.
\]
Le premier opérateur compare des secteurs séparés par trois positions ; le
second, par quatre. Les noms de canal 90 et 120 se réfèrent ici à
ces séparations dans le cycle orienté, et non à une identification avec un
opérateur physique d'un autre domaine.

\begin{proposition}[Complétude du système transversal]
\label{prop:k-transversal}
L'opérateur
\[
 \mathcal D=(D_3,D_4,Q):\mathbb Z^{12}\longrightarrow\mathbb Z^{25}
\]
a un rang égal à douze et des facteurs invariants non nuls
\[
 \operatorname{diag}(1,\ldots,1,12),
\]
avec onze unités. En particulier, leurs sorties déterminent un unique registre dodécaphasique.
\end{proposition}

\begin{proof}
Si \(D_3z=D_4z=0\), le vecteur est constant le long des deux pas.
Comme \(3\) et \(4\) engendrent le cycle modulo douze,
\(\ker D_3\cap\ker D_4=\langle\mathbf1\rangle\). La charge totale
élimine cette dernière liberté parce que \(Q(\mathbf1)=12\).

Pour l'affirmation entière, les lignes de \((D_3,D_4)\) sont des incidences orientées d'un graphe connexe. Un arbre couvrant fournit onze facteurs invariants unitaires. Tout mineur d'ordre douze doit contenir la ligne \(Q\). En ajoutant les onze premières colonnes à la dernière, celle-ci s'annule dans les lignes d'incidence et vaut douze dans la ligne \(Q\). Tous les mineurs maximaux sont divisibles par douze et l'arbre en produit un de module exactement douze. Le dernier facteur est donc douze.
\end{proof}

La vérification sur \eqref{eq:k-vector} donne
\[
 D_3K=b^{90},\qquad D_4K=b^{120},\qquad Q(K)=6263.
\]
Les formules \eqref{eq:k-sumas-armonicas} et
\eqref{eq:k-bloques-firmados} reconstruisent à leur tour \(U\) à partir de ces
différences. On ferme ainsi un triangle de représentations exactes du
même objet : registre transversal, vecteur signé et registre dodécaphasique. L'égalité
des résultats ne fait pas de leurs codomaines un seul et même domaine ; elle démontre que
les applications déclarées restituent l'information pertinente.

\subsection{Deux lectures rationnelles : \texorpdfstring{\(\kappa_{36}\)}{kappa36} et \texorpdfstring{\(\kappa_{\mathrm{per}}\)}{kappa périodique}}
\label{subsec:k-kappas}

Une fois \(K\) construit, la coordinatisation en base mille permet deux opérations
distinctes. La première lit un seul tour ; la seconde prolonge périodiquement le registre
dodécaphasique. Soit \(B=1000\) et définissons son entier de concaténation
\begin{equation}
 N_K=\sum_{m=1}^{12}K_mB^{12-m}
 =234543140729659824621058914794146601.
 \label{eq:k-numerador}
\end{equation}
Les zéros initiaux d'un bloc, comme celui de 058, font partie de son écriture
à trois chiffres. Ils ne changent pas sa valeur entière, mais permettent de concaténer
sans ambiguïté les douze blocs.

\begin{definition}[Lectures finie et périodique du registre dodécaphasique]
\label{def:k-kappas}
La lecture d'un tour est
\[
 \kappa_{36}=\sum_{m=1}^{12}K_mB^{-m}=\frac{N_K}{B^{12}}.
\]
Le prolongement périodique est
\[
 K_j^{\mathrm{per}}=K_{1+((j-1)\bmod12)},\qquad
 \kappa_{\mathrm{per}}=\sum_{j\geq1}K_j^{\mathrm{per}}B^{-j}.
\]
\end{definition}

\begin{proposition}[Valeur et période exactes du registre dodécaphasique prolongé]
\label{prop:k-periodo}
On a
\begin{equation}
 \kappa_{\mathrm{per}}=\frac{N_K}{B^{12}-1},\qquad
 \kappa_{\mathrm{per}}-\kappa_{36}
 =\frac{N_K}{B^{12}(B^{12}-1)}
 =B^{-12}\kappa_{\mathrm{per}}.
 \label{eq:k-diferencia-kappas}
\end{equation}
Le mot infini du registre dodécaphasique a pour période minimale douze en base mille et pour période minimale trente-six en base dix.
\end{proposition}

\begin{proof}
La somme du premier tour est \(N_KB^{-12}\) ; chaque nouveau tour
multiplie sa contribution par \(B^{-12}\). La série géométrique produit
\(N_K/(B^{12}-1)\), et la soustraction donne la seconde identité.

Les douze composantes de \(K\) sont distinctes. Une période propre du cycle
identifierait deux d'entre elles ; la période minimale des blocs est donc
douze. Le développement décimal est canonique, car le registre dodécaphasique n'a pas de queue
de blocs 999. Si sa période décimale minimale était \(d\), elle diviserait
36. Un regroupement par trois produirait une période de blocs
\(d/\gcd(d,3)\). Aucun diviseur propre de 36 ne donne douze par cette
opération. La période décimale minimale est donc 36.
\end{proof}

Les deux lectures sont rationnelles, mais ne sont pas le même rationnel. La première se termine après douze blocs ; la seconde répète le tour. En outre, \(0<\kappa_{\mathrm{per}}-\kappa_{36}<B^{-12}\). Leur proximité ne permet pas de substituer l'une à l'autre dans une identité exacte. En particulier, la clôture finie d'alpha utilise \(\kappa_{36}\), tandis que la somme de la section prolongée utilise \(\kappa_{\mathrm{per}}\).

% Fragmento conservado; véase MANIFIESTO_ANTECEDENTES_LECTURA.json.
\subsection{Réversibilité de la constante rationnelle et changement de phase}

La constante rationnelle \(\kappa_{\mathrm{per}}\) détermine intégralement
le registre \(K\) lorsque la base, la longueur et l'origine
de lecture sont conservées. Sa fonction numérique et sa fonction structurelle sont reliées
par une inversion exacte.

\begin{proposition}[Récupération entière et transformation cyclique]
Soit \(I(K)=\sum_{j=1}^{12}K_j1000^{12-j}\). Alors
\[
 I(K)=(1000^{12}-1)\kappa_{\mathrm{per}}.
\]
Le développement entier de \(I(K)\) en douze blocs de base \(1000\)
récupère tous les \(K_j\), y compris les zéros initiaux de chaque bloc.
Si \(\rho(K)=(K_2,\ldots,K_{12},K_1)\), on a
\begin{equation}
 \kappa(\rho K)=1000\kappa(K)-K_1.
 \label{eq:k-rotacion-racional}
\end{equation}
\end{proposition}
\begin{proof}
La première identité inverse la définition de l'évaluation périodique.
La division euclidienne itérée par \(1000\), avec douze positions fixées,
récupère le vecteur. D'autre part,
\[
 I(\rho K)=1000I(K)-K_1(1000^{12}-1).
\]
La division par \(1000^{12}-1\) donne \eqref{eq:k-rotacion-racional}.
\end{proof}

La récupération compose donc \(\kappa\mapsto K\mapsto U\)
et les transformations \(D_3K,D_4K,Q\) déjà construites. L'évaluation
rationnelle conserve le registre signé et les différences cycliques qui
en dépendent. La profondeur totale et les opérateurs de transport d'une
histoire restent dans leurs coordonnées propres : l'inversion précédente
a pour domaine le registre entier dodécaphasique.

Écrivons \(\kappa_j=\kappa(\rho^jK)\), avec des indices modulo \(12\).
La même identité donne
\[
 K_{j+1}=1000\kappa_j-\kappa_{j+1},\qquad
 \sum_{j=0}^{11}\kappa_j=\frac{\sum_{j=1}^{12}K_j}{999}
 =\frac{6263}{999}.
\]
La phase agit de manière exacte sur la constante : sa valeur est fixe
dans la lecture canonique, tandis que ses rotations suivent une loi affine.
La somme de l'orbite est invariante sous le choix de l'origine.

L'incidence ultérieure possède également une expression dans ces
coordonnées. Le seuil cubique détermine
\[
 B_K=\{j+1:1000\kappa_j-\kappa_{j+1}\ge729\}
 =\{4,6,9,10\}.
\]
La tétrade se récupère à partir de l'orbite rationnelle parce que celle-ci
récupère d'abord les composantes entières. Le support signé associé
à \(\alpha\) utilise en outre ses registres régionaux et la normalisation
de frontière. La relation entre coordonnée et incidence s'exprime ainsi
au moyen d'opérations identifiables.

C'est une définition opérationnelle de la fonction de \(K\) : il fournit
une mise en coordonnées entières réversible du registre orienté, dont
l'évaluation périodique et les fonctions d'incidence participent aux
constructions ultérieures. La valeur, la transformation et le domaine
sont conservés conjointement.



\subsection{Périodicité et conoyau de l'opérateur de transport}
\label{subsec:k-clase-acarreo}

La périodicité que nous venons de démontrer appartient au registre dodécaphasique. Elle n'implique
pas la périodicité des coordonnées régionales, des retenues de la clôture
ni d'alpha. Il est utile de préciser cette séparation au moyen d'un second
quotient, purement entier.

Soit \(S\) le décalage cyclique de douze coordonnées, d'orientation fixée, et soit \(b\geq2\) une base entière. L'opérateur de retenue périodique est \(\partial_b=S-bI\). Si une identité périodique s'écrit
\[
 A=P+E-\Phi-K+\partial_bC,
\]
alors les transformations
\[
 K\longmapsto K+\partial_bh,\qquad C\longmapsto C+h
\]
conservent son membre de droite. Cette invariance compare les représentants d'une
classe de retenue ; elle n'autorise pas à modifier le registre dodécaphasique déjà
sélectionné par le registre.

\begin{proposition}[Quotient de la retenue périodique]
\label{prop:k-cociente-acarreo}
Avec l'orientation précédente,
\[
 \operatorname{coker}(S-bI)\cong\mathbb Z/(b^{12}-1)\mathbb Z.
\]
\end{proposition}

\begin{proof}
Le module cyclique se présente sous la forme
\(\mathbb Z[t]/(t^{12}-1)\). Imposer \(S=bI\) équivaut à imposer
\(t=b\) et laisse la relation \(b^{12}-1=0\). L'évaluation des
coefficients en \(b\), avec l'ordre de concaténation choisi, représente
la classe résultante.
\end{proof}

Le dénominateur \(B^{12}-1\) de \(\kappa_{\mathrm{per}}\) et ce quotient proviennent de la même longueur cyclique, mais leurs objets restent distincts : l'un est une évaluation rationnelle et l'autre un quotient d'un module de retenues. La clôture finie d'alpha sera une chaîne orientée à frontière droite. Son prolongement transportera une frontière distante compatible, et non une retenue identifiée par décret entre la première et la treizième position.
